% TOP_PROBLEM: {"id": 2307018, "title": "Research Problems in Function Theory — Problem 7.18", "category_id": 8, "category": {"id": 8, "name": "analysis", "display_name": "Analysis", "description": "Limits, continuity, calculus, and function theory.", "slug": "analysis", "order_index": 8, "created_at": "2026-07-31T15:26:25.671Z"}, "status": "open"}
% FIRST_POSTED: null
% ENTRY_KIND: statement_only
% Imported from: batch05.tex; UnsolvedMath ID 2307018
\documentclass[11pt]{article}
\usepackage{amsmath,amssymb,mathtools}
\usepackage{fontspec,unicode-math}
\setmainfont{Latin Modern Roman}
\setmathfont{Latin Modern Math}
\usepackage{xurl,hyperref}
\newcommand{\sref}[2]{\hyperlink{source:#1:#2}{[S#2]}}
\newcommand{\cataloguescope}[1]{\paragraph{Mathematical statement.}#1}
\begin{document}
% UnsolvedMath ID 2307018; source catalogue code AMR-022-7018
\setcounter{section}{2307017}
\section{Research Problems in Function Theory — Problem 7.18}\label{2307018}
{\small\sffamily UnsolvedMath ID 2307018 · Analysis}
\subsection{Definitions and mathematical statement}

\paragraph{Shared notation.}$\mathbb N=\{1,2,\ldots\}$, and $\mathbb Z,\mathbb Q,\mathbb R,\mathbb C$ denote the integers, rationals, reals and complex numbers. Any other conventions are specified locally.


Let $\Gamma\subset\mathbb C$ be a Jordan curve, meaning the image of a continuous injective parametrisation of the circle, with $0$ in its bounded complementary component. Write $C(\Gamma)$ for the complex continuous functions with the uniform norm. Since $0\notin\Gamma$, the functions $z\mapsto z^n$ are continuous on $\Gamma$ for every $n\in\mathbb Z$. A family spans $C(\Gamma)$ when its finite complex linear combinations are dense in this norm.

\paragraph{Statement recorded in the supplied source.}
Characterise the curves $\Gamma$ for which there exists an infinite set $E\subset\mathbb Z$ such that
\[\overline{\operatorname{span}\{z^n|_\Gamma:n\in\mathbb Z\setminus E\}}^{\|\cdot\|_\infty}=C(\Gamma).\]
In other words, under what conditions on $\Gamma$ can infinitely many Laurent powers be omitted without losing density? \sref{2307018}{2}
The original notes Wermer's density result after deleting the constant power for curves of infinite length, and related results for finitely many deleted powers. These are context for the infinite-deletion question. \sref{2307018}{2}
\subsection{Short English statement}
Determine which Jordan curves permit deleting infinitely many integer powers from a uniformly dense Laurent-power family.
\subsection{Sources}
\begin{description}
\item[\hypertarget{source:2307018:1}{S1}] ULAM.AI and the UnsolvedMath Contributors, \emph{UnsolvedMath: A Curated Collection of Open Mathematics Problems}, record 2307018 (AMR-022-7018). CC BY4.0. \url{https://huggingface.co/datasets/ulamai/UnsolvedMath}.
\item[\hypertarget{source:2307018:2}{S2}] W. K. Hayman and E. F. Lingham, \emph{Research Problems in Function Theory} (2018), Chapter7, Problem7.18 and its complete update. \url{https://arxiv.org/abs/1809.07200}.
\end{description}
\label{end:2307018}
\end{document}
